\section{TYPE-002: as any Assertions in Model Handling}
\label{sec:type-002}
\subsection*{Metadata}
\begin{tabular}{ll}
\textbf{Issue ID} & TYPE-002 \\
\textbf{Severity} & MEDIUM \\
\textbf{Category} & Type Safety \\
\textbf{Branch} & \branchref{fix/TYPE-002-remove-as-any-in-model-handling}
\textbf{Changes} & \branchdiff{fix/TYPE-002-remove-as-any-in-model-handling} \\
\end{tabular}

\subsection*{Executive Summary}
At \srcfile{src/loader.ts:418}, the model's \code{provider} property is forcibly set via \code{(model as any).provider = "openai"}. This bypasses the TypeScript type system and mutates the \code{Model} object in-place, violating the type contract. The \code{as any} is unnecessary — the \code{provider} field already allows \code{string} assignment. If \code{Model} is refactored to use a getter or readonly property, this silent mutation could cause bugs.

\subsection*{Root Cause Analysis}
The \code{as any} cast is used unnecessarily:

\begin{lstlisting}[language=TypeScript]
// BEFORE — unnecessary as any cast
(model as any).provider = "openai";
\end{lstlisting}

\subsection*{Fix Implementation}
The \code{as any} cast is removed:

\begin{lstlisting}[language=TypeScript]
// AFTER — direct property assignment
model.provider = "openai";
\end{lstlisting}

\subsection*{Affected Files}
\begin{itemize}
    \item \srcfile{src/loader.ts} — Line 418
    \item \srcfile{src/\_\_tests\_\_/type-002.test.ts}
\end{itemize}

\subsection*{Verification}
TypeScript compilation passes without the \code{as any} cast.
